\section{AI w badaniach klinicznych}

% SLIDE ID: sec09-goal
\BlockGoalFrame{\SlideID{sec09-goal}}{
  \item zobaczyć realne, a nie marketingowe zastosowania \texttt{AI} w badaniach klinicznych,
  \item zrozumieć, gdzie \texttt{AI} może wspierać operacje i jakość danych,
  \item uporządkować ograniczenia, walidację i odpowiedzialność człowieka.
}

% SLIDE ID: sec09-what-is-ai
\begin{frame}{Czym AI może być w badaniach klinicznych}
  \SlideID{sec09-what-is-ai}
  \begin{itemize}
    \item zestaw metod wspierających rozpoznawanie wzorców, klasyfikację i analizę treści,
    \item narzędzie do pracy na danych operacyjnych, klinicznych i dokumentacyjnych,
    \item system podpowiadający ryzyko, a nie magiczny moduł podejmujący decyzje za zespół,
    \item warto myśleć o \texttt{AI} przez pryzmat konkretnego procesu, nie przez modne etykiety.
  \end{itemize}
\end{frame}

% SLIDE ID: sec09-ml-nlp-genai
\begin{frame}{Machine Learning, NLP, Generative AI}
  \SlideID{sec09-ml-nlp-genai}
  \begin{itemize}
    \item \texttt{Machine Learning}: wykrywanie wzorców, scoring ryzyka, predykcje.
    \item \texttt{NLP}: analiza tekstu, dokumentów, narracji medycznych i korespondencji.
    \item \texttt{Generative AI}: szkice podsumowań, sugestie treści, wyszukiwanie w dokumentach.
    \item W praktyce najważniejsze pytanie brzmi: \emph{co system rekomenduje i kto to zatwierdza?}
  \end{itemize}
\end{frame}

% SLIDE ID: sec09-ai-assistant
\begin{frame}{AI jako asystent, nie zastępca}
  \SlideID{sec09-ai-assistant}
  \begin{itemize}
    \item \texttt{AI} może wskazać problem, priorytet albo możliwe wyjaśnienie.
    \item Nie powinna samodzielnie zamykać \texttt{query}, klasyfikować \texttt{SAE} ani podejmować decyzji medycznej.
    \item Człowiek pozostaje właścicielem decyzji, dokumentacji i odpowiedzialności regulacyjnej.
    \item Dobra automatyzacja skraca pracę zespołu, ale nie usuwa nadzoru.
  \end{itemize}
\end{frame}

% SLIDE ID: sec09-where-ai-works
\begin{frame}{Gdzie AI działa dziś}
  \SlideID{sec09-where-ai-works}
  \begin{itemize}
    \item \texttt{data cleaning} i wykrywanie anomalii w danych,
    \item automatyczne sugestie \texttt{query} dla niespójnych wpisów,
    \item \texttt{risk-based monitoring} i priorytetyzacja ośrodków,
    \item predykcja rekrutacji i analiza opóźnień,
    \item analiza dokumentów, medical coding i analiza obrazów.
  \end{itemize}
\end{frame}

% SLIDE ID: sec09-data-cleaning
\begin{frame}{Data cleaning i wykrywanie anomalii}
  \SlideID{sec09-data-cleaning}
  \begin{itemize}
    \item modele mogą wykrywać nietypowe zakresy, brakujące wzorce i niespójne sekwencje danych,
    \item system może porównywać zachowanie ośrodka z resztą badania,
    \item to wsparcie dla data managera, a nie automatyczne stwierdzenie błędu,
    \item wynik modelu powinien prowadzić do pytania, nie do bezrefleksyjnej akcji.
  \end{itemize}
\end{frame}

% SLIDE ID: sec09-query-rbm
\begin{frame}{Query suggestions i risk-based monitoring}
  \SlideID{sec09-query-rbm}
  \begin{itemize}
    \item \texttt{AI} może proponować treść \texttt{query} na podstawie reguł i podobnych przypadków.
    \item Może też wskazywać ośrodki o podwyższonym ryzyku operacyjnym lub jakościowym.
    \item W \texttt{RBM} przydatne są sygnały o opóźnieniach, brakach danych i nietypowych odchyleniach.
    \item Nadal potrzebna jest decyzja człowieka: czy to realny sygnał, czy fałszywy alarm.
  \end{itemize}
\end{frame}

% SLIDE ID: sec09-recruitment-selection
\begin{frame}{Rekrutacja, selekcja pacjentów i dokumenty}
  \SlideID{sec09-recruitment-selection}
  \begin{itemize}
    \item predykcja rekrutacji pomaga ocenić tempo i ryzyko niedowiezienia targetu,
    \item algorytmy mogą wspierać wyszukiwanie potencjalnie pasujących pacjentów,
    \item \texttt{NLP} pomaga przeszukiwać dokumenty, protokoły i kryteria włączenia / wyłączenia,
    \item tu szczególnie ważne są jakość źródeł, bias i kontrola nad kryteriami.
  \end{itemize}
\end{frame}

% SLIDE ID: sec09-coding-imaging
\begin{frame}{Medical coding i analiza obrazów}
  \SlideID{sec09-coding-imaging}
  \begin{itemize}
    \item \texttt{AI} może wspierać mapowanie terminów do słowników medycznych,
    \item może też pomagać w klasyfikacji treści opisowych i priorytetyzacji review,
    \item w badaniach obrazowych wspiera analizę wzorców i wstępne oznaczanie zmian,
    \item obszary te są użyteczne, ale wymagają szczególnie mocnej walidacji i kontroli jakości.
  \end{itemize}
\end{frame}

% SLIDE ID: sec09-risky
\begin{frame}{Gdzie AI jest ryzykowne}
  \SlideID{sec09-risky}
  \begin{itemize}
    \item tam, gdzie wynik modelu jest nieczytelny i trudny do uzasadnienia,
    \item tam, gdzie dane uczące są niepełne, stronnicze albo historycznie złej jakości,
    \item tam, gdzie system zaczyna wpływać na decyzje kliniczne bez jasnego nadzoru,
    \item tam, gdzie użytkownik ufa sugestii bardziej niż danym źródłowym.
  \end{itemize}
\end{frame}

% SLIDE ID: sec09-validation
\begin{frame}{Co trzeba walidować}
  \SlideID{sec09-validation}
  \begin{itemize}
    \item dane wejściowe i ich jakość,
    \item logikę modelu lub zasady generowania odpowiedzi,
    \item powtarzalność wyników i granice działania,
    \item sposób prezentacji rekomendacji użytkownikowi,
    \item ślad decyzji: kto zaakceptował, odrzucił lub zmienił sugestię.
  \end{itemize}
\end{frame}

% SLIDE ID: sec09-explainability-bias
\begin{frame}{Explainability, bias i zgodność regulacyjna}
  \SlideID{sec09-explainability-bias}
  \begin{itemize}
    \item trzeba rozumieć, na jakiej podstawie system wskazuje ryzyko lub rekomendację,
    \item trzeba ocenić, czy model nie wzmacnia istniejących błędów i nierówności danych,
    \item trzeba zapewnić dokumentację, wersjonowanie i kontrolę zmian,
    \item zgodność regulacyjna oznacza możliwość obrony procesu podczas audytu lub inspekcji.
  \end{itemize}
\end{frame}

% SLIDE ID: sec09-human-responsibility
\begin{frame}{Odpowiedzialność człowieka pozostaje kluczowa}
  \SlideID{sec09-human-responsibility}
  \begin{itemize}
    \item \texttt{AI} może rekomendować, ale to człowiek odpowiada za decyzję.
    \item Zespół musi wiedzieć, kiedy zaufać sugestii, a kiedy ją odrzucić.
    \item Procedury powinny jasno opisywać użycie, review i eskalację wyjątków.
    \item Bez tego nawet dobry model staje się ryzykiem operacyjnym i regulacyjnym.
  \end{itemize}
\end{frame}

% SLIDE ID: sec09-demo
\DemoFrame{AI w operacjach badania klinicznego}{
  \begin{itemize}
    \item przykład wykrywania nietypowych wzorców danych,
    \item przykład sugestii \texttt{query} dla niespójnego wpisu,
    \item przykład dashboardu ryzyka dla \texttt{RBM},
    \item przykład wyszukiwania informacji w dokumentach badania.
  \end{itemize}
}

% SLIDE ID: sec09-case
\CaseStudyFrame{AI wykrywa nietypowy wzorzec danych w ośrodku}{
  Model wskazuje, że jeden ośrodek ma wyjątkowo mało odchyleń laboratoryjnych, bardzo szybkie uzupełnianie formularzy i nietypowo jednorodne dane między pacjentami.
}{
  Czy to sygnał jakościowy, problem operacyjny, a może cecha populacji? Jakie działania powinien podjąć zespół zanim wyciągnie wnioski lub uruchomi eskalację?
}

% SLIDE ID: sec09-wrap
\begin{frame}{Najważniejsza myśl z bloku AI}
  \SlideID{sec09-wrap}
  \begin{itemize}
    \item Najbardziej użyteczne \texttt{AI} w badaniach klinicznych wspiera konkretne procesy.
    \item Największe ryzyko pojawia się wtedy, gdy zespół nie rozumie ograniczeń modelu.
    \item Dobre wdrożenie \texttt{AI} wymaga walidacji, przejrzystości i stałego nadzoru człowieka.
  \end{itemize}
\end{frame}
